% Ad Familiares 7.13 --- headnote
% ad-familiares-07-13, 4 March 53 BC, Rome

\textit{Cicero to C. Trebatius Testa in Gaul, written from Rome on
4 March 53 BC. Trebatius has evidently complained that Cicero's
correspondence had lapsed and has read the silence as pique;
Cicero answers that the lapse was practical (he did not know
where his friend was), then turns the complaint back on its
author. The little flight of mock indignation in the first
paragraph --- ``which is it that makes you so haughty, your purse,
or that the imperator consults you?'' --- belongs to the running
joke of the Trebatius series: Cicero pretending to take umbrage
at the airs his young protégé has acquired on Caesar's staff, and
playing on the double sense of \textit{consuli}, ``to be consulted''
as a jurist and ``to be looked after'' as a friend, with the
extra pun on \textit{inaurari}, ``to be gilded.''}

\textit{The second paragraph is sustained legal joking. In Gaul,
Cicero says, property is not claimed by the formal civil
procedure of \textit{manum consertum} --- the symbolic laying-on of
hands at the disputed object --- but by the sword; and the standard
interdict against forcible entry, with its exception ``provided
that you have not been the first to come with armed men and
force,'' is no protection where the disputants are barbarians.
The closing pun is on the Treveri tribe (whose territory lay near
Caesar's winter quarters) and on the old legal tag \textit{aere,
argento, auro}, ``of bronze, silver, and gold,'' from the formulae
about coined metal: the Treveri, he is told, are \textit{capital}
fellows --- meaning both deadly and, in legal parlance, liable to
the capital sentence --- but he would rather they were the harmless
kind associated with the formula. The lawyerly wit is the whole
point of the letter, and the matter-of-fact dating subscription
\textit{D. IIII non. Mart.} is preserved as Cicero wrote it.}
